%! Modeling with Quadratics — Unit 2, Lesson 2.7 — Exit Ticket
\documentclass[10pt]{article}
\usepackage{algebra2-article}
\usepackage{algebra2-boxes}
\newcommand{\TallMath}[1]{$\displaystyle #1\rule[-1.4em]{0pt}{3.2em}$}

\begin{document}

\pageheader{Unit 2, Lesson 2.7}{Exit Ticket}

\vspace{0.12in}

No notes. Show your steps. Answer the actual question --- with units.

\begin{enumerate}[leftmargin=1.8em, itemsep=14pt]
  \item \textbf{Projectile.} A ball is thrown up at $64$ ft/s from an $80$-ft roof: $h(t)=-16t^2+64t+80$.
        \begin{enumerate}[label=(\alph*), leftmargin=1.9em, itemsep=4pt]
          \item What is the \emph{maximum height}, and when does it occur? \hfill \blank{4.2cm}
          \item When does the ball \emph{hit the ground}? \hfill \blank{3.2cm}
        \end{enumerate}

  \item \textbf{Maximum area.} A rectangular pen uses $100$ ft of fence on all four sides, with width $x$ and
        length $50-x$, so $A(x)=x(50-x)$. What width gives the \emph{maximum area}, and what is that area?
        \[ x=\blank{2.2cm} \qquad\qquad \text{max area}=\blank{3.0cm} \]

  \item \textbf{SOL-style multiple choice.} For the projectile $h(t)=-16t^2+80t+96$, which quantity gives the
        \emph{maximum height} the object reaches?
        \begin{enumerate}[label=\Alph*., leftmargin=1.9em, itemsep=1pt]
          \item $h(0)$
          \item the $t$-coordinate of the vertex
          \item the $h$-coordinate of the vertex
          \item the positive zero of $h(t)$
        \end{enumerate}
\end{enumerate}

\end{document}
